%%%%%%%%%%%%%%%%%%%%%%%%%%%%%%%%%%%%%%%%%%%%%%%%%%
%% Praca magisterska — Politechnika Poznańska   %%
%% Wydział Informatyki i Telekomunikacji        %%
%% Autorka: Maria Nowicka (151851)              %%
%% Promotor: prof. dr hab. inż. Robert Wrembel  %%
%% Temat: Odkrywanie zależności między          %%
%%        obiektami w bazie danych              %%
%%%%%%%%%%%%%%%%%%%%%%%%%%%%%%%%%%%%%%%%%%%%%%%%%%

\documentclass[polish,master,a4paper,oneside]{ppfcmthesis}

\usepackage[utf8]{inputenc}
\usepackage[OT4]{fontenc}
\usepackage{csquotes}

% ── Dodatkowe pakiety ────────────────────────────────────────
\usepackage{amsmath}
\usepackage{amssymb}
\usepackage{booktabs}
\usepackage{tabularx}
\usepackage{multirow}
\usepackage{longtable}
\usepackage{xcolor}
\usepackage{listings}
\usepackage{float}

\lstset{
  basicstyle=\ttfamily\small,
  keywordstyle=\color{blue},
  commentstyle=\color{gray},
  stringstyle=\color{teal},
  breaklines=true,
  frame=single,
  numbers=left,
  numberstyle=\tiny\color{gray},
  captionpos=b
}

% ── Strona tytułowa ─────────────────────────────────────────
\author{Maria Nowicka \album{151851}}
\authortitle{}

\title{Odkrywanie zależności między obiektami w~bazie danych}

\ppsupervisor{prof.~dr hab.~inż.~Robert Wrembel}

\ppyear{2026}

% ============================================================
\begin{document}

% ── Strona tytułowa ─────────────────────────────────────────
\frontmatter\pagestyle{empty}%
\maketitle\cleardoublepage%

% ── Miejsce na kartę pracy dyplomowej ───────────────────────
\thispagestyle{empty}\vspace*{\fill}%
\begin{center}
  Tutaj będzie karta pracy dyplomowej;\\
  oryginał wstawiamy do wersji dla archiwum PP,\\
  w pozostałych kopiach wstawiamy ksero.
\end{center}%
\vfill\cleardoublepage%

% ── Streszczenie (PL) ───────────────────────────────────────
\chapter*{Streszczenie}
\addcontentsline{toc}{chapter}{Streszczenie}

% UWAGA REDAKCYJNA: Streszczenie zawiera cel, problem, lukę badawczą i metodykę.
%   Brakuje wyników liczbowych — uzupełnić po zakończeniu eksperymentów.
%   Streszczenie nie zawiera cytowań (standard akademicki).
%   Docelowa długość: 250–300 słów. Obecna wersja: ok. 220 słów.

% TODO (WYMAGANE): Uzupełnić ostatni akapit o konkretne wartości F1/AUC-ROC.

Data lineage i~data provenance to~pojęcia opisujące śledzenie przepływu
danych i~historię transformacji danych w~systemach bazodanowych. Grafy
lineage reprezentują obiekty bazodanowe — tabele, widoki, kolumny — jako
węzły, a~transformacje danych jako krawędzie. Znajomość zależności między
obiektami w~bazie danych jest kluczowa z~punktu widzenia zapewnienia
jakości danych, audytu procesów ekstrakcji, transformacji i~ładowania
danych (ang. \emph{Extract, Transform, Load}, ETL), debugowania potoków
danych oraz spełnienia wymogów regulacyjnych takich jak RODO.

Głównym problemem rozważanym w~niniejszej pracy jest zjawisko tzw.\
\emph{broken lineage} — zerwania ciągłości grafu zależności. Do~zerwania
ciągłości dochodzi, gdy w~procesach przetwarzania danych używane są tabele
tymczasowe lub funkcje definiowane przez użytkownika (ang.
\emph{User-Defined Function}, UDF). Tabele tymczasowe nie są rejestrowane
w~standardowych słownikach bazy danych, a~UDF ukrywają wewnętrzną logikę
transformacji. W~obu przypadkach automatyczne odtworzenie pełnego grafu
zależności nie jest możliwe — w~grafie lineage pojawiają się brakujące
krawędzie, których nie można ustalić na~podstawie dostępnych metadanych.

Istniejące narzędzia klasy data governance umożliwiają śledzenie lineage
w~czasie rzeczywistym, lecz wymagają dostępu do~pełnych metadanych
i~nie rozwiązują problemu brakujących krawędzi. W~literaturze naukowej
podejmowane są próby rozwiązania tego problemu przy użyciu analizy logów
systemowych, przepisywania zapytań SQL oraz uczenia maszynowego.
Dotychczasowe podejścia nie obejmują jednak analizy wpływu struktury grafu
lineage na~dobór algorytmów predykcji. Grafy lineage mają charakter
dwudzielny — krawędzie przepływu danych łączą wyłącznie węzły tabela
i~zadanie ETL — co powoduje, że~klasyczne heurystyki grafowe oparte
na~wspólnych sąsiadach są w~tym kontekście nieskuteczne.

W~ramach niniejszej pracy zaimplementowano i~porównano trzy grupy
algorytmów predykcji brakujących krawędzi w~grafach lineage: heurystyki
grafowe (Common Neighbors, Jaccard, Adamic-Adar, Preferential Attachment,
L3), algorytmy uczenia maszynowego oparte na~cechach topologicznych grafu
(Random Forest, RUSBoost, LightGBM) oraz metody zanurzeń sieciowych
(\emph{network embeddings}) reprezentowane przez Node2Vec. Tę ostatnią
grupę celowo odróżniono od~grafowych sieci neuronowych (GNN), z~którymi
bywa~mylona — Node2Vec uczy reprezentacji wektorowych węzłów poprzez
biased random walks, bez~mechanizmu propagacji wiadomości
(\emph{message passing}) charakterystycznego dla GNN. Algorytmy
poddano ocenie z~użyciem metryk F1, precyzji, czułości oraz AUC-ROC
i~AUC-PR na~zbiorze danych DLG-DG-23, zawierającym 18~grafów lineage
ze~środowiska Huawei Cloud.

% TODO (WYMAGANE): Dodać zdanie z wynikami, np.:
%   „Najlepsze wyniki uzyskał algorytm X, osiągając wartość F1 równą Y
%    w scenariuszu Z."

\bigskip
\noindent\textbf{Słowa kluczowe:} data lineage, data provenance, broken
lineage, predykcja krawędzi, uczenie maszynowe, grafy dwudzielne,
Node2Vec, RUSBoost.

% ── Abstract (EN) ───────────────────────────────────────────
\chapter*{Abstract}
\addcontentsline{toc}{chapter}{Abstract}

% TODO (WYMAGANE): Add a sentence with numerical results after experiments,
%   e.g., "The best-performing algorithm achieved an F1 score of X in scenario Y."

Data lineage and data provenance refer to tracking data flow and
transformation history in database systems. Lineage graphs represent
database objects — tables, views, columns — as nodes and data
transformations as edges. Knowledge of dependencies between database
objects is essential for data quality assurance, auditing
Extract-Transform-Load (ETL) processes, debugging data pipelines, and
compliance with regulations such as GDPR.

The main problem addressed in this thesis is the phenomenon of
\emph{broken lineage} — gaps in the continuity of the dependency graph.
These gaps occur when data processing pipelines use temporary tables or
user-defined functions (UDFs). Temporary tables are not registered in
standard database catalogs, and UDFs hide the internal transformation
logic. In both cases, automatic reconstruction of the full dependency
graph is not possible — the lineage graph contains missing edges that
cannot be determined from available metadata.

Existing data governance tools support real-time lineage tracking, but
they require access to complete metadata and do not address the problem
of missing edges. Research literature attempts to solve this problem using
system log analysis, SQL query rewriting, and machine learning. However,
existing approaches do not analyze the impact of lineage graph structure
on algorithm selection. Lineage graphs have a bipartite structure —
data-flow edges connect only table--job node pairs — which makes classical
neighborhood-based graph heuristics ineffective in this context.

In this thesis, three groups of algorithms for predicting missing edges in
lineage graphs are implemented and compared: graph heuristics (Common
Neighbors, Jaccard, Adamic-Adar, Preferential Attachment, L3), machine
learning algorithms based on topological graph features (Random Forest,
RUSBoost, LightGBM), and a network embedding method represented by
Node2Vec. The latter group is deliberately distinguished from
Graph Neural Networks (GNNs) with which it is sometimes confused —
Node2Vec learns node representations via biased random walks, without
the message-passing mechanism characteristic of GNNs. The algorithms
are evaluated using F1, precision, recall, AUC-ROC, and AUC-PR
metrics on the DLG-DG-23 dataset, comprising 18~lineage graphs from
a~Huawei Cloud environment.

\bigskip
\noindent\textbf{Keywords:} data lineage, data provenance, broken lineage,
link prediction, machine learning, bipartite graphs, Node2Vec, RUSBoost.

% ── Spis treści ─────────────────────────────────────────────
\pagenumbering{Roman}\pagestyle{ppfcmthesis}%
\tableofcontents*
\cleardoublepage

% ============================================================
%  ROZDZIAŁY GŁÓWNE
% ============================================================
\mainmatter

% ============================================================
%  ROZDZIAŁ 1 — WSTĘP
% ============================================================
\chapter{Wstęp}

Współczesne systemy bazodanowe przetwarzają dane za~pośrednictwem złożonych
potoków transformacji SQL. W~środowiskach korporacyjnych dane są wielokrotnie
agregowane i~przekształcane przez sekwencje zapytań, a~wyniki pośrednie
przechowywane w~tabelach tymczasowych. Tymczasowość tych obiektów sprawia,
że~po~zakończeniu procesu przetwarzania ślad po~wykonanych transformacjach
zanika — niemożliwe staje się ustalenie, skąd pochodzi dana wartość
w~tabeli wynikowej ani jakie operacje ją~wytworzyły. Zjawisko to, określane
w~literaturze jako \emph{data lineage} bądź \emph{data provenance}, stanowi
istotny problem zarówno z~perspektywy audytu danych, jak i~ich
jakości~\cite{Psallidas2020,Yamada2023}.

Formalnie problem ten można wyrazić w~kategoriach teorii grafów: graf
zależności (ang. \emph{lineage graph}) reprezentuje przepływ danych między
obiektami bazy — tabelami, widokami, zapytaniami i~ich atrybutami.
Gdy~pewne węzły lub~krawędzie zostaną utracone (np.~z~powodu usunięcia
tabeli tymczasowej lub zastosowania funkcji definiowanej przez użytkownika),
mówi się o~zjawisku \emph{broken lineage}~\cite{Boinski2025}. Celem
niniejszej pracy jest zastosowanie algorytmów predykcji brakujących krawędzi
w~celu odtworzenia przerwanych połączeń w~grafie lineage.

% TODO: Rozważyć dodanie krótkiego przykładu motywującego (1–2 zdania
%       z~konkretnym scenariuszem biznesowym, np. raportowanie finansowe).

\section{Cel i~zakres pracy}

Celem pracy jest zbadanie skuteczności istniejących algorytmów predykcji
brakujących krawędzi w~grafach lineage. Zakres obejmuje trzy klasy metod:

\begin{itemize}
  \item heurystyki grafowe (ang. \emph{graph heuristics}),
  \item klasyczne metody uczenia maszynowego (ang. \emph{Machine Learning}, ML),
  \item embeddingi węzłów grafu (ang. \emph{node embeddings}).
\end{itemize}

Eksperymenty przeprowadzono na~zbiorze danych DLG-DG-23, obejmującym
18~zanonimizowanych grafów dwudzielnych ze~środowiska Huawei Cloud
\cite{Chen2023}. Przyjęto trzy scenariusze symulacji brakujących krawędzi,
odzwierciedlające typowe przypadki broken lineage:

\begin{enumerate}
  \item losowe usunięcie 20\% krawędzi ze~zbioru treningowego,
  \item symulację ukrycia zależności przez tabelę pośrednią (staging),
  \item symulację ukrycia krawędzi do~tabeli docelowej (data mart).
\end{enumerate}

Zakres pracy nie~obejmuje opracowania autorskiego algorytmu predykcji,
odtworzenia rzeczywistego schematu bazy danych ani analizy treści
zanonimizowanych danych.

% TODO: Jeśli scenariusze mają ustalone nazwy (np. Scenariusz A, B, C)
%       w rozdziale 3 — ujednolicić nazewnictwo w całej pracy.

\section{Struktura pracy}

Praca składa się z~pięciu rozdziałów. Rozdział~2 zawiera przegląd pojęć
podstawowych — data lineage, data provenance, broken lineage oraz predykcja
krawędzi — a~także omówienie wybranych prac z~literatury i~stosowanych
metryk ewaluacji. Rozdział~3 opisuje problem badawczy, architekturę
zaproponowanego rozwiązania, metodykę symulacji brakujących krawędzi oraz
zaimplementowane algorytmy. Rozdział~4 prezentuje wyniki eksperymentów
przeprowadzonych na~zbiorze DLG-DG-23 wraz z~ich analizą. Rozdział~5
zawiera wnioski końcowe, podsumowanie uzyskanych wyników oraz propozycje
dalszych kierunków badań.

% ============================================================
%  ROZDZIAŁ 2 — PODSTAWY TEORETYCZNE I PRZEGLĄD LITERATURY
% ============================================================
\chapter{Podstawy teoretyczne i~przegląd literatury}

Niniejszy rozdział wprowadza pojęcia podstawowe niezbędne do~zrozumienia
dalszych części pracy. W~sekcji~\ref{sec:lineage_def} zdefiniowano pojęcia
data lineage i~data provenance. Sekcja~\ref{sec:graf_lineage} opisuje
strukturę grafów lineage na~przykładzie zbioru danych DLG-DG-23.
Sekcja~\ref{sec:broken_lineage} omawia przyczyny niekompletności grafów,
a~sekcja~\ref{sec:link_prediction} przedstawia problem predykcji krawędzi.
Sekcja~\ref{sec:metryki} opisuje metryki ewaluacji, a~sekcja~\ref{sec:literatura}
zawiera przegląd wybranej literatury.

\section{Data lineage i~data provenance}
\label{sec:lineage_def}

Pojęcia \emph{data lineage} i~\emph{data provenance} są w~literaturze
stosowane zarówno zamiennie, jak i~z~rozróżnieniem semantycznym~\cite{Psallidas2020}.
W~węższym ujęciu \emph{data lineage} odnosi się do~śledzenia przepływu
informacji przez kolejne etapy przetwarzania — opisuje, przez jakie
transformacje przeszły dane i~jak zmieniała się ich postać. Z~kolei
\emph{data provenance} kładzie nacisk na~pochodzenie i~autentyczność danych:
rejestruje, skąd dane zostały pobrane i~kto je~wytworzył. W~praktyce
granica między tymi pojęciami jest płynna i~większość opracowań traktuje je
jako synonimy~\cite{Yamada2023}. Takie podejście przyjęto również w~niniejszej
pracy.

Znajomość historii danych ma~istotne znaczenie praktyczne. W~systemach
raportowych możliwość odtworzenia ścieżki transformacji pozwala zlokalizować
źródło błędnych wartości — bez~tej informacji ręczna weryfikacja wielkich
zbiorów danych jest nieefektywna. Innym przykładem jest bezpieczne usuwanie
obiektów bazy: tabela pozornie niepowiązana z~żadnym procesem może
w~rzeczywistości dostarczać danych do~raportów końcowych, a~jej usunięcie
spowoduje utratę danych bez~widocznego komunikatu~błędu~\cite{Boinski2025}.

% TODO: Można rozszerzyć o zastosowania w zgodności z regulacjami (GDPR,
%       audyt finansowy) — jeśli autorka chce poszerzyć kontekst biznesowy.

\section{Struktura grafu lineage w~zbiorze DLG-DG-23}
\label{sec:graf_lineage}

Zbiór danych DLG-DG-23, opracowany przez Chen et~al.~\cite{Chen2023},
zawiera 18~grafów lineage pobranych ze~środowiska produkcyjnego Huawei
Cloud. Dane zostały zanonimizowane — nie~udostępniono schematów tabel.
Grafy różnią się skalą: liczba węzłów waha się od~278 do~17~085.

Każdy graf zawiera trzy typy węzłów:

\begin{itemize}
  \item \textbf{job} — zapytanie SQL lub~zadanie ETL (ang. \emph{Extract,
        Transform, Load}),
  \item \textbf{table} — tabela lub~widok przechowujący dane,
  \item \textbf{field} — atrybut (kolumna) należący do~tabeli.
\end{itemize}

Krawędzie grafu przyjmują dwa typy:

\begin{itemize}
  \item \textbf{data\_flow} — krawędź skierowana między tabelą a~jobem
        lub~w~kierunku odwrotnym; reprezentuje przepływ danych,
  \item \textbf{parent\_child} — krawędź między tabelą a~jej atrybutem
        (polem); wyraża relację przynależności kolumny do~tabeli.
\end{itemize}

Graf jest dwudzielny w~tej~części, w~której węzły job i~table tworzą
rozłączne zbiory połączone wyłącznie krawędziami data\_flow. Właściwość
dwudzielności odzwierciedla strukturę przetwarzania ETL: dane przepływają
między magazynami danych (table) a~procesami transformacji (job).

% TODO: Warto dodać rysunek poglądowy — przykładowy mały graf z węzłami
%       job, table, field i krawędziami data_flow, parent_child.
%       Rysunek musi być wprowadzony zdaniem w tekście przed wystąpieniem.

% TODO: Dodać tabelę ze statystykami 18 grafów (liczba węzłów, krawędzi,
%       gęstość). Tabela musi być poprzedzona zdaniem „Tabela X przedstawia...".

\section{Niekompletność grafów lineage — broken lineage}
\label{sec:broken_lineage}

Graf lineage bywa niekompletny z~kilku powodów. Tabele tymczasowe,
tworzone na~potrzeby pojedynczej sesji lub~procesu ETL, znikają po~jego
zakończeniu. Jeśli pośredniczyły w~przepływie danych między dwoma trwałymi
tabelami, znika wraz z~nimi odpowiedni fragment grafu — krawędź
reprezentująca tę~zależność nigdy nie~zostaje utrwalona w~metadanych.
Analogiczny skutek mają widoki zmaterializowane tworzone tymczasowo oraz
funkcje definiowane przez użytkownika (ang. \emph{User-Defined Function},
UDF), które opakowują logikę transformacji bez~jawnego zapisu przepływu
danych~\cite{Yamada2023}.

Zjawisko to, określane jako \emph{broken lineage}, prowadzi do~sytuacji,
w~której graf zawiera węzły bez~połączeń lub~z~niepełnym sąsiedztwem,
mimo że~rzeczywista zależność między odpowiadającymi im~obiektami istnieje.
W~niniejszej pracy problem broken lineage jest modelowany jako zadanie
predykcji brakujących krawędzi — każdej usuniętej krawędzi odpowiada para
węzłów, dla~której algorytm ma~ocenić prawdopodobieństwo istnienia
połączenia~\cite{Boinski2025}.

% TODO: Rozważyć dodanie formalnej definicji problemu:
%       G = (V, E) — kompletny graf, G' = (V, E \ E_b) — graf z broken lineage,
%       cel: dla każdej pary (u,v) \in V \times V \setminus E' oszacować P((u,v) \in E).

\section{Predykcja brakujących krawędzi}
\label{sec:link_prediction}

Predykcja brakujących krawędzi (ang. \emph{link prediction}) to~zadanie
uczenia maszynowego polegające na~ocenie prawdopodobieństwa istnienia
krawędzi między parą węzłów w~grafie~\cite{KohanMarzagao2023}. W~klasycznym
ujęciu problem ten pojawia się w~grafach, w~których część krawędzi nigdy
nie~została zaobserwowana. W~niniejszej pracy przyjęto odmienne podejście:
dysponując pełnymi grafami lineage, krawędzie są usuwane w~sposób
kontrolowany zgodnie z~przyjętymi scenariuszami symulacji (rozdział~3),
a~zadaniem algorytmów jest ich odtworzenie.

W~pracy zastosowano trzy klasy algorytmów:

\begin{enumerate}
  \item \textbf{Heurystyki grafowe} — miary podobieństwa sąsiedztwa węzłów,
        takie jak Common Neighbors, indeks Jaccarda i~miara Adamic-Adar;
        nie~wymagają uczenia, opierają się wyłącznie na~strukturze grafu.
  \item \textbf{Klasyczne metody ML} — Random Forest, RUSBoost i~LightGBM;
        wykorzystują wektory cech opisujących pary węzłów,
        w~tym cechy strukturalne grafu i~metryki podobieństwa nazw
        (np.~Jaro-Winkler).
  \item \textbf{Zanurzenia sieciowe} (ang. \emph{network embeddings})
        — metoda Node2Vec, która odwzorowuje węzły grafu w~przestrzeni
        wektorowej tak, by~węzły o~podobnym kontekście strukturalnym były
        reprezentowane przez zbliżone wektory; predykcja krawędzi odbywa się
        przez klasyfikator (MLP) operujący na~iloczynie Hadamarda
        embeddingów pary węzłów. Należy podkreślić, że~Node2Vec nie~jest
        grafową siecią neuronową (GNN) — uczy reprezentacji wyłącznie
        poprzez biased random walks, bez~mechanizmu propagacji wiadomości
        (\emph{message passing}). Rozszerzenie pracy o~właściwą rodzinę GNN
        (np.~R-GCN, HGT) wskazano w~rozdziale podsumowującym jako kierunek
        dalszych badań.
\end{enumerate}

% TODO: Opisać dokładniej każdy algorytm w podrozdziałach (2.4.1–2.4.3)
%       lub przenieść szczegółowy opis do rozdziału 3.

\section{Metryki ewaluacji}
\label{sec:metryki}

Do~oceny jakości predykcji krawędzi stosuje się standardowe metryki
klasyfikacji binarnej. W~niniejszej pracy przyjęto następujące metryki:

\begin{itemize}
  \item \textbf{Precision} (precyzja) — frakcja poprawnie przewidzianych
        krawędzi spośród wszystkich krawędzi wskazanych przez model,
  \item \textbf{Recall} (czułość) — frakcja poprawnie wykrytych krawędzi
        spośród wszystkich rzeczywiście brakujących krawędzi,
  \item \textbf{F1} — średnia harmoniczna precyzji i~czułości; główna
        miara porównawcza stosowana w~pracy,
  \item \textbf{AUC-ROC} — pole pod krzywą ROC; mierzy zdolność modelu
        do~rozróżnienia par połączonych od~niepołączonych.
\end{itemize}

Metryka F1 jest szczególnie istotna w~warunkach niezrównoważonych klas,
charakterystycznych dla~grafów rzadkich, gdzie liczba potencjalnych
brakujących krawędzi wielokrotnie przewyższa liczbę krawędzi
rzeczywistych~\cite{Boinski2025,KohanMarzagao2023}.

% TODO: Sprawdzić, czy AUC-ROC jest faktycznie raportowane w kodzie —
%       jeśli nie, usunąć lub zastąpić metryką AUC-PR.

\section{Przegląd literatury}
\label{sec:literatura}

Niniejsza sekcja omawia siedem prac wybranych w~ramach systematycznego
przeglądu literatury (ang. \emph{Systematic Literature Review}, SLR).

% TODO: Rozbudować każdy akapit do ok. 3–5 zdań.

Boiński et~al.~\cite{Boinski2025} zaproponowali zastosowanie klasycznych
metod ML do~problemu broken lineage. Autorzy zastosowali algorytm RUSBoost
(ang. \emph{Random Under-Sampling Boosting}) na~cechach obejmujących
metadane tabel oraz miary podobieństwa nazw (Jaro-Winkler), uzyskując
wartość F1 równą 0{,}79. Praca ta~stanowi bezpośrednie odniesienie
metodologiczne dla~niniejszej pracy.

Chen et~al.~\cite{Chen2023} udostępnili zbiór danych DLG-DG-23 — 18~grafów
lineage pobranych ze~środowiska produkcyjnego Huawei Cloud, obejmujących
grafy o~liczbie węzłów od~278 do~17~085. Zbiór ten~stanowi podstawowy
materiał eksperymentalny niniejszej pracy.

Dietrich et~al.~\cite{Dietrich2023} omówili specyfikę śledzenia danych w~SQL
z~rekurencyjnymi wspólnymi wyrażeniami tabelowymi (ang. \emph{Common Table
Expression}, CTE). Autorzy zaproponowali technikę przepisywania zapytań
(ang. \emph{query rewriting}) umożliwiającą propagację informacji
o~proweniencji przez operacje rekurencyjne.

Kohan~Marzagão et~al.~\cite{KohanMarzagao2023} zaproponowali jądro grafowe
(ang. \emph{Provenance Graph Kernel}, PGK) do~porównywania grafów
proweniencji. Podejście to~umożliwia pomiar podobieństwa między grafami
lineage bez~konieczności ich jawnego wyrównywania.

Psallidas et~al.~\cite{Psallidas2020} opracowali system OneProvenance,
który ekstrahuje informacje o~proweniencji danych bezpośrednio z~logów
silnika bazy danych, budując strukturę drzewa zapytań (ang. \emph{QQTree})
bez~modyfikacji kodu źródłowego systemu.

Shan et~al.~\cite{Shan2022} zaproponowali model KPI-HGNN oparty
na~heterogenicznej grafowej sieci neuronowej (ang. \emph{Heterogeneous
Graph Neural Network}, HGNN). Model stosuje mechanizm wykrywania
społeczności (ang. \emph{community detection}) w~celu agregacji informacji
o~lokalnym sąsiedztwie węzłów.

Yamada et~al.~\cite{Yamada2023} rozszerzyli standardowe podejścia do~data
lineage o~obsługę funkcji definiowanych przez użytkownika (UDF). Autorzy
zaproponowali mechanizm rozszerzonych adnotacji lineage (ang. \emph{augmented
lineage}), który pozwala śledzić przepływ danych przez ciała funkcji UDF
niewidocznych w~standardowych planach zapytań.

% ============================================================
%  ROZDZIAŁ 3 — PROBLEM I SPOSÓB JEGO ROZWIĄZANIA
% ============================================================
\chapter{Problem i~sposób jego rozwiązania}

Niniejszy rozdział opisuje architekturę zaproponowanego rozwiązania,
pochodzenie i~charakterystykę danych, metodykę symulacji zjawiska broken
lineage oraz szczegóły implementacji zrealizowanych algorytmów.

\section{Charakterystyka danych}
\label{sec:dane}

% TODO: Opisać zbiór DLG-DG-23 szczegółowo: statystyki poszczególnych grafów
%       (tabela z nazwami grafów, liczbą węzłów, krawędzi, gęstością).
%       Wprowadzić tabelę zdaniem: „Tabela X przedstawia...".

\section{Scenariusze symulacji broken lineage}
\label{sec:scenariusze}

% TODO: Opisać trzy scenariusze:
%       (1) losowe usunięcie 20% krawędzi (split 80/20),
%       (2) Scenariusz A — usunięcie tabeli staging,
%       (3) Scenariusz B/C — ukrycie krawędzi do data martu / widoku.
%       Dla każdego: motywacja, sposób implementacji, liczba usuniętych krawędzi.

\section{Zaimplementowane algorytmy}
\label{sec:algorytmy}

% TODO: Opisać każdy algorytm w osobnym podrozdziale (3.3.1–3.3.3).
%       Wymagana struktura: (a) opis metody, (b) wektor cech lub parametry,
%       (c) hiperparametry użyte w eksperymentach.

\subsection{Heurystyki grafowe}

% TODO: Common Neighbors, Jaccard, Adamic-Adar — wzory matematyczne,
%       interpretacja, złożoność obliczeniowa.

\subsection{Klasyczne metody uczenia maszynowego}

% TODO: Random Forest, RUSBoost, LightGBM — opis, cechy wejściowe (tabela),
%       obsługa niezrównoważonych klas.

\subsection{Zanurzenia sieciowe — Node2Vec}

% TODO: Opis metody Node2Vec (biased random walks, parametry p i q),
%       generowanie embeddingu krawędzi (iloczyn Hadamarda),
%       klasyfikator MLP. Wyjaśnić różnicę względem GNN
%       (brak message passing, uczenie na sekwencjach a nie agregacji
%       sąsiedztwa).

\section{Architektura rozwiązania}
\label{sec:architektura}

% TODO: Schemat blokowy całego potoku: wczytanie grafu →
%       symulacja broken lineage → ekstrakcja cech → trening →
%       predykcja → ewaluacja.
%       Rysunek musi być wprowadzony w tekście przed jego wystąpieniem.

% ============================================================
%  ROZDZIAŁ 4 — OCENA ROZWIĄZANIA
% ============================================================
\chapter{Ocena rozwiązania}

Niniejszy rozdział przedstawia wyniki eksperymentów przeprowadzonych
zgodnie z~metodologią opisaną w~rozdziale~3. Dla~każdego z~trzech
scenariuszy symulacji oraz każdego z~trzech algorytmów podano wartości
metryk F1, Precision, Recall i~AUC-ROC uśrednione po~wszystkich
18~grafach zbioru DLG-DG-23.

\section{Wyniki eksperymentów}
\label{sec:wyniki}

% TODO: Uzupełnić wynikami po zakończeniu eksperymentów.
%       Tabela zbiorcza: jeden wiersz na algorytm × scenariusz.
%       Wprowadzić tabelę zdaniem w tekście.

\section{Analiza wyników}
\label{sec:analiza}

% TODO: Porównanie algorytmów: który działa najlepiej dla jakiego scenariusza.
%       Odwołanie do literatury (porównanie z F1=0.79 z \cite{Boinski2025}).
%       Omówienie przypadków trudnych — grafy małe vs. duże.

% ============================================================
%  ROZDZIAŁ 5 — UWAGI KOŃCOWE
% ============================================================
\chapter{Uwagi końcowe}

Niniejszy rozdział zawiera podsumowanie uzyskanych wyników, wnioski
dotyczące skuteczności zastosowanych algorytmów oraz propozycje
dalszych kierunków badań.

\section{Podsumowanie}
\label{sec:podsumowanie}

% TODO: Odpowiedzieć na pytania badawcze PB1–PB4 sformułowane we wstępie.
%       Co zadziałało, co nie zadziałało i dlaczego.

\section{Kierunki dalszych badań}
\label{sec:dalsze_badania}

% TODO: Propozycje rozszerzeń:
%       - autorski algorytm dedykowany grafom dwudzielnym,
%       - uwzględnienie schematów tabel (jeśli dane staną się dostępne),
%       - GNN / GraphSAGE jako następny krok po Node2Vec,
%       - ewaluacja na innych zbiorach danych.

% ============================================================
%  LITERATURA
% ============================================================
\bibliographystyle{plain}{\raggedright\sloppy\small\bibliography{bibliography}}

% ============================================================
%  STRONA PRAW AUTORSKICH (wymagana przez szablon PP)
% ============================================================
\ppcolophon

\end{document}
